\begin{abstract}
L'alternating Criteria Search è un algorimto euristico che permette di ottenere
soluzioni di alta qualità per problemi generici di programmazione lineare intera
mista.
Ad oggi le implementazioni di ottimizzatori per problemi MIP allo stato dell'arte,
sfruttano algoritmi euristici come strumento per migliorare l'efficienza nella
ricerca di una soluzione ottima.
L'ACS può essere utilizzato in congiunzione con un algoritmo di ricerca ottimale,
oppure come risolutore per quella classe di problemi in
cui il raggiungimento dell'ottimo risulta essere computazionalmente troppo oneroso.

L'algoritmo presenta molti parametri che devono essere impostati: 
i valori assegnati e la strategia con la quale questi vengono individuati
influenzano le prestazioni dell'algoritmo.
L'obiettivo è quello di ricavare, per via sperimentale, la combinazione di
valori che fornisce le migliori prestazioni.
\end{abstract}
